% Chapter 7 -- reading sites and scoring. Update with web/* and scoring/* (semantic_index.py), gui/results_table.py,
% gui/detail_panel.py, gui/table_sort.py, gui/context_menu.py, gui/geo_filter.py, gui/search_box.py,
% gui/social_links.py, gui/lazy_map.py (map view of Results), companies.workplace.
\chapter{Reading sites and scoring}\label{ch:scoring}

\section{Reading a company website}
Sites are read on several threads (\emph{companies in parallel}, 1--16, default 8),
with the rules of \code{web/fetch.py}: public http/https hosts only, robots.txt
respected, a delay between requests to the same domain (default 2\,s), at most 5 pages
per site, pages cached 7 days, a domain that answers HTTP 429 is not asked again.
Engines: \emph{local} (default, httpx); \emph{playwright} and \emph{firecrawl}
(only with a key; only the URL is sent) are tried when a page looks JavaScript-only.

\begin{enumerate}[leftmargin=1.5em]
  \item \textbf{Careers page}: from link texts and paths, then common paths, then the
        sitemap declared in robots.txt.
  \item \textbf{Postings cascade}: ATS board $\rightarrow$ schema.org JobPosting
        (JSON-LD) $\rightarrow$ heuristics on the careers page $\rightarrow$ an
        ``open application'' signal only.
  \item \textbf{ATS boards} read through their public APIs: Greenhouse, Lever, Ashby,
        Workable, Recruitee, SmartRecruiters, Personio, Workday, BambooHR. When a
        careers page has no ATS link, boards named like the domain are probed; a guess
        counts only if the board's data names the site's own domain.
  \item \textbf{Posting attributes}: contract, work mode, salary (only when published,
        never estimated), business travel.
\end{enumerate}

\section{Requirements: must, nice, gate}
Posting text is split into units and classified by heading (``Requirements'' $\to$
must, ``Nice to have'' $\to$ nice, duties, offer) and cue words. Each unit gets an
anchored level with the CV line quoted: 4 strong, 3 clear, 2 partial, 1 weak,
0 explicit contradiction, \emph{none} = no evidence (uncertain, never a failure).
The \textbf{blocking gate} fires only when a must-have of a hard kind (language,
driving licence, own car, business travel) is \emph{explicitly contradicted} by the
CV. Silence never blocks.

\section{The scores in Results}
\begin{tabularx}{\textwidth}{@{}l Y@{}}
\toprule
\hdr \thead{Column} & \thead{Meaning (0--100 unless noted)} \\
\midrule
Hiring       & Hiring / Maybe / No signal / Website unreachable / No website. \\
Match        & Sorting value. With postings: 0.50 fit + 0.25 desirability + 0.25 relevance (reduced if blocked). Without: discounted mix of spontaneous, relevance, desirability. Scaled by confidence. \\
Fit          & Best posting vs CV: title and skills, requirement coverage, factors for unmet requirements. \\
Desirability & Your preferences: distance, contract, work mode, salary, travel. Unknown = neutral. \\
Relevance    & Company vs profile: semantic similarity of site text, sector, CV words, ESCO role coverage. \\
Salary, Contract, Work mode & From the best posting. \\
Distance, Job ads & Kilometres from the start point; number of postings. \\
Confidence   & 0--1: how much was actually read (text, pages, posting source). \\
Affinity, Verdict, Fits as, Type & How well the CV fits what the company needs, its verdict, the best direction and the kind of company (chapter~\ref{ch:affinity}); empty until a CV is loaded. \\
Revenue trend & Yearly growth of revenue from the stored figures, e.g.\ ``+6.3\,\%/year (2021--2024)'' (section~\ref{sec:trend}). \\
\bottomrule
\end{tabularx}

\vspace{3pt}
\textbf{Table or map.} A button switches table and map (section~\ref{sec:map}); the detail
panel lists postings with their \emph{Match with my CV} badge (section~\ref{sec:match}, not
the company \emph{Match} above) and \ui{Simulate screening} (chapter~\ref{ch:analysis}).

Filters: name, hiring only, website only, hiring first, remote/hybrid, permanent, salary. The
detail panel shows reasons with their source, gaps and the job ads, a \textbf{row of links}
beside the company name (section~\ref{sec:links}), and the buttons \ui{Open website},
\ui{Open careers page}, \ui{Open job ad}, \ui{Open LinkedIn} and \ui{Write e-mail\ldots}
(chapter~\ref{ch:apply}). \ui{Open LinkedIn} only opens a search in \emph{your} browser.

\textbf{Workplace.} CSV export, \ui{My companies} and favourites have a \emph{Workplace}
column: the best posting's location, else the company's town, else blank.

\section{Working with the tables}\label{sec:tables}
\textbf{Sorting.} Click a column header: ascending, a second click descending, with an arrow.
Numbers sort as numbers (``9\,\%'' before ``10\,\%'', a range by its lower bound), dates as dates,
text without regard to case or accents, and empty cells always last. The order survives refills
and is remembered per table until the program closes; there is no filter row. With no column
chosen, Results, \ui{My companies} and the favourites are ordered by affinity
(section~\ref{sec:order}).

\textbf{Right-click menus.} Every table has one, built from the same actions as the buttons and
the menu bar (a disabled action is left out): Results (website, careers, LinkedIn, e-mail, ad,
simulate screening), \ui{My companies} (create full card, e-mail, find similar, save, delete),
the favourites, \ui{Previous searches}, the job search and the interest table. The Delete key runs
the delete action where there is one.

\textbf{Place selectors.} Above \ui{My companies}, the favourites, Results, the job-search
results and \ui{Previous searches} sit three linked multi-choice menus, \ui{Country},
\ui{Region} and \ui{Province}, each with counts (``Emilia-Romagna (12)'') and an \emph{Unknown
place} entry. Region and Province exist for Italy only and are enabled when Italy is chosen;
about 50 countries are known, a foreign place keeps just country and town as written, and
nothing uses the network. The selectors combine (AND) with the column filters and the search box.

\textbf{Search by meaning, everywhere.} The same search box is in all those lists. Typing runs
an exact-word pass at once and, after 0.4\,s, a pass \emph{by meaning} in the background with the
local model: a hybrid ranking. A whole phrase found as words (or its Italian/English variant)
scores 75--100\,\%, all the words in one passage 60--85\,\%, meaning alone at most 59\,\% (a hit
under 35\,\% is hidden), so an exact hit never ranks below an unrelated neighbour. Texts are cut
into sentence-like passages (about 280 characters, at most 60 per document) and every vector is
cached in the store: a repeated search only reads the query, a changed card only its changed
passages. A hit shows its relevance and the \emph{matching sentence} with where it was found, for
example ``87\,\% $\cdot$ Presentation: <the sentence>'' (a \ui{Matched text} column in \ui{My
companies}, the row tooltip in Results). Without the model only exact words are found and the box
says so.

\begin{warnbox}
Without the scoring model the search runs with a built-in lexical embedder: confidence is
lowered and a ``Reduced accuracy'' notice with \ui{Why?} appears (chapter~\ref{ch:troubleshooting}).
\end{warnbox}
